\documentclass[11pt]{article}
\usepackage[margin=1in]{geometry}
\usepackage{amsmath,amssymb,amsthm}
\usepackage{algorithm}
\usepackage{algpseudocode}
\usepackage{booktabs}
\usepackage{hyperref}

\newtheorem{definition}{Definition}
\newtheorem{theorem}{Theorem}
\newtheorem{lemma}[theorem]{Lemma}
\newtheorem{property}{Property}

\title{\textbf{Mathematical Specification, Gradient Dynamics, and Invariance Theorems of Overlap Segmentation Losses (Soft Dice and Soft Jaccard/IoU) (ALGO-NN-21)}}
\author{Team Scaibu \\ \small Email: \texttt{jaydeep.9664920749@gmail.com} \\ \small Architecture and Core Engine Team}
\date{\today}

\begin{document}
\maketitle

\begin{abstract}
Semantic and volumetric medical segmentation often suffers from overwhelming background class dominance, where standard cross-entropy fails to prioritize small foreground lesions. Continuous overlap losses directly optimize set-theoretic similarity metrics (S{\o}rensen-Dice and Jaccard Index). This specification provides the formal continuous relaxation, derives analytic gradients, proves background invariance, provides line-by-line pseudocode matching the reference implementation, and establishes complexity bounds.
\end{abstract}

\section{Notation and Mathematical Preliminaries}
\label{sec:notation}

Let $\mathbf{p} \in [0, 1]^N$ denote predicted foreground probabilities and $\mathbf{g} \in \{0, 1\}^N$ ground-truth binary indicators across $N$ pixels or voxels. Let $\epsilon \in \mathbb{R}_{>0}$ denote the Laplace smoothing term.

\subsection{Continuous Set Intersection and Union}
\paragraph{Intuitive Meaning:} Rather than thresholding probabilities into discrete $0/1$ decisions, Soft Dice and Soft IoU multiply continuous probabilities with binary ground truth to compute a differentiable intersection $|\mathbf{p} \cap \mathbf{g}|$.

\paragraph{Formal Mathematical Definition:}
\begin{itemize}
    \item \textbf{Soft Dice Loss:}
    \begin{equation}
    \label{eq:soft_dice}
    \mathcal{L}_{\mathrm{Dice}}(\mathbf{p}, \mathbf{g}; \epsilon) = 1 - \frac{2 \sum_{i=1}^N p_i g_i + \epsilon}{\sum_{i=1}^N p_i + \sum_{i=1}^N g_i + \epsilon} = 1 - \frac{2 |\mathbf{p} \cap \mathbf{g}|_\epsilon}{|\mathbf{p}|_1 + |\mathbf{g}|_1 + \epsilon}.
    \end{equation}
    \item \textbf{Soft Jaccard / IoU Loss:}
    \begin{equation}
    \label{eq:soft_iou}
    \mathcal{L}_{\mathrm{IoU}}(\mathbf{p}, \mathbf{g}; \epsilon) = 1 - \frac{\sum_{i=1}^N p_i g_i + \epsilon}{\sum_{i=1}^N p_i + \sum_{i=1}^N g_i - \sum_{i=1}^N p_i g_i + \epsilon}.
    \end{equation}
\end{itemize}

\subsection{Summary of Symbols}
\begin{itemize}
    \item $N \in \mathbb{N}$: Total number of flattened pixels or voxels.
    \item $\mathbf{p} \in [0, 1]^N$: Predicted probability vector.
    \item $\mathbf{g} \in [0, 1]^N$: Ground truth target mask.
    \item $\epsilon \in \mathbb{R}_{>0}$: Additive Laplace smoothing constant.
    \item $\mathrm{DSC} \in [0, 1]$: Continuous S{\o}rensen-Dice coefficient.
    \item $\mathrm{IoU} \in [0, 1]$: Continuous Jaccard index.
    \item $\mathbf{G} = \nabla_{\mathbf{p}} \mathcal{L} \in \mathbb{R}^N$: Analytic gradient vector.
\end{itemize}

\section{Problem Definition and Core Concepts}
\label{sec:problem_definition}

\subsection{Analytic Gradient Derivation}
Differentiating $\mathcal{L}_{\mathrm{Dice}}$ with respect to probability $p_k$:
Let $U = 2 \sum p_i g_i + \epsilon$ and $V = \sum p_i + \sum g_i + \epsilon$:
\begin{equation}
\frac{\partial \mathcal{L}_{\mathrm{Dice}}}{\partial p_k} = -\frac{2 g_k V - U}{V^2} = -\frac{2 g_k \left(\sum p_i + \sum g_i + \epsilon\right) - \left(2 \sum p_i g_i + \epsilon\right)}{\left(\sum p_i + \sum g_i + \epsilon\right)^2}.
\end{equation}

\section{Algorithmic Specification}
\label{sec:algorithm}

Algorithm~\ref{alg:overlap_losses} specifies the complete calculation procedure matching \texttt{impl.py}.

\begin{algorithm}[H]
\caption{Continuous Overlap Segmentation Loss and Gradient Evaluation}
\label{alg:overlap_losses}
\begin{algorithmic}[1]
\Require Predicted probabilities $\mathbf{p} \in [0, 1]^N$, targets $\mathbf{g} \in [0, 1]^N$, loss type $\mathrm{type} \in \{\text{dice}, \text{iou}\}$, smoothing $\epsilon > 0$.
\Ensure Loss $\mathcal{L}$, Dice score $\mathrm{DSC}$, IoU score $\mathrm{IoU}$, gradients $\mathbf{G} \in \mathbb{R}^N$.
\State $I \gets \sum_{i=1}^N p_i g_i, \quad S_p \gets \sum_{i=1}^N p_i, \quad S_g \gets \sum_{i=1}^N g_i$
\State $D_{\mathrm{num}} \gets 2.0 \cdot I + \epsilon, \quad D_{\mathrm{den}} \gets S_p + S_g + \epsilon$
\State $\mathrm{DSC} \gets \frac{D_{\mathrm{num}}}{D_{\mathrm{den}}}, \quad \mathcal{L}_{\mathrm{Dice}} \gets 1.0 - \mathrm{DSC}$
\State $J_{\mathrm{num}} \gets I + \epsilon, \quad J_{\mathrm{den}} \gets S_p + S_g - I + \epsilon$
\State $\mathrm{IoU} \gets \frac{J_{\mathrm{num}}}{J_{\mathrm{den}}}, \quad \mathcal{L}_{\mathrm{IoU}} \gets 1.0 - \mathrm{IoU}$
\State Initialize $\mathbf{G} \gets []$
\If{$\mathrm{type} = \text{dice}$}
    \State $\mathcal{L} \gets \mathcal{L}_{\mathrm{Dice}}$
    \For{$i = 1$ to $N$}
        \State $g_{\mathrm{val}} \gets -\frac{2.0 \cdot g_i \cdot D_{\mathrm{den}} - D_{\mathrm{num}}}{D_{\mathrm{den}}^2}$
        \State Append $g_{\mathrm{val}}$ to $\mathbf{G}$
    \EndFor
\Else
    \State $\mathcal{L} \gets \mathcal{L}_{\mathrm{IoU}}$
    \For{$i = 1$ to $N$}
        \State $g_{\mathrm{val}} \gets -\frac{g_i \cdot J_{\mathrm{den}} - (1.0 - g_i) \cdot J_{\mathrm{num}}}{J_{\mathrm{den}}^2}$
        \State Append $g_{\mathrm{val}}$ to $\mathbf{G}$
    \EndFor
\EndIf
\State \Return $\mathcal{L}, \quad \mathrm{DSC}, \quad \mathrm{IoU}, \quad \mathbf{G}$
\end{algorithmic}
\end{algorithm}

\section{Theoretical Guarantees and Background Invariance}
\label{sec:guarantees}

\begin{theorem}[Invariance to Pure Background Regions]
\label{thm:dice_background_invariance}
\textbf{Assumptions:} Let true background pixels satisfy $g_j = 0, p_j = 0$. Suppose the background canvas is expanded by $M$ empty pixels ($g_{N+k} = 0, p_{N+k} = 0$).
\textbf{Guarantee:} The loss $\mathcal{L}_{\mathrm{Dice}}$ and foreground gradients $\frac{\partial \mathcal{L}_{\mathrm{Dice}}}{\partial p_i}$ are strictly invariant to canvas size $M$.
\textbf{Proof:}
For any empty background pixel $k$, $p_k g_k = 0, p_k = 0, g_k = 0$.
Thus $\sum_{i=1}^{N+M} p_i g_i = \sum_{i=1}^N p_i g_i$ and $\sum_{i=1}^{N+M} (p_i + g_i) = \sum_{i=1}^N (p_i + g_i)$.
The numerator $D_{\mathrm{num}}$ and denominator $D_{\mathrm{den}}$ remain unchanged, preserving both the scalar loss and all component gradients identically.
\textbf{Limitations:} Soft Dice is non-convex; when initial predictions are all near zero, gradients on foreground pixels can be small, requiring combination with cross-entropy.
\end{theorem}

\section{Complexity Analysis}
\label{sec:complexity}

\begin{itemize}
    \item \textbf{Time Complexity:} $\mathcal{O}(N)$ sequential floating-point multiply-accumulate operations.
    \item \textbf{Space Complexity:} $\mathcal{O}(N)$ gradient buffer storage.
\end{itemize}

\section{Worked Numerical Example}
\label{sec:worked_example}

Let $\mathbf{p} = [0.8, 0.2], \mathbf{g} = [1.0, 0.0], \epsilon = 1.0$:
\begin{enumerate}
    \item $I = 0.8(1.0) + 0.2(0.0) = 0.8$.
    \item $S_p = 0.8 + 0.2 = 1.0, S_g = 1.0 + 0.0 = 1.0$.
    \item $D_{\mathrm{num}} = 2(0.8) + 1.0 = 2.6, D_{\mathrm{den}} = 1.0 + 1.0 + 1.0 = 3.0$.
    \item $\mathrm{DSC} = \frac{2.6}{3.0} = 0.866667 \implies \mathcal{L}_{\mathrm{Dice}} = 1 - 0.866667 = 0.133333$.
    \item Gradients:
    $g_1 = -\frac{2(1.0)(3.0) - 2.6}{3.0^2} = -\frac{6.0 - 2.6}{9.0} = -\frac{3.4}{9.0} = -0.377778$.
    $g_2 = -\frac{2(0.0)(3.0) - 2.6}{9.0} = +\frac{2.6}{9.0} = +0.288889$.
\end{enumerate}

\section{Numerical Considerations and Edge Cases}
\label{sec:numerical}

\begin{itemize}
    \item \textbf{Smoothing Denominator Epsilon:} $\epsilon = 1.0$ ensures that empty target masks ($S_g = 0$) with zero predictions ($S_p = 0$) evaluate to $\mathrm{DSC} = \frac{0 + 1}{0 + 1} = 1.0 \implies \mathcal{L} = 0.0$.
\end{itemize}

\begin{thebibliography}{9}
\bibitem{milletari2016} F.~Milletari, N.~Navab, and S.-A.~Ahmadi, ``V-Net: Fully Convolutional Neural Networks for Volumetric Medical Image Segmentation,'' \emph{3DV}, 2016.
\end{thebibliography}

\end{document}
